% SPDX-License-Identifier: Apache-2.0
\section{\code{txhdl::types}}
\label{sec:r-types}

Values, and the two markers. \code{Bit} is two valued, \code{Logic} is
IEEE 1164's nine, and \code{U<N>} and \code{I<N>} are the numeric
vectors, with the width a const parameter. A width that is an
expression of other widths, \code{U<\{A + B\}>}, needs nightly Rust,
which is the largest constraint the embedding found, and the
prototype declares each width by hand instead. The conversions from
Rust integers are what let a literal be written where a value is
expected, and the narrowing ones are checked in debug builds.

\code{Transaction} marks a struct that moves over a channel, and a
bare word is one too. \code{Tag} is a synchronisation domain and its
policy, as a type; which clock edge moves the data is a
\code{Clock}, in the next module.

\lstinputlisting[caption={\code{lib/src/types.rs}.},label={lst:r-types}]{lib/src/types.rs}

\section{\code{txhdl::netlist}}
\label{sec:r-netlist}

The lowering's target, Section~\ref{sec:r-lowering}: \code{Expr},
\code{Target} and \code{Stmt}, which \code{\#[lower]} builds; the
\code{Lowered} unit that holds them with its fields, ports, wires and
memory contents; the two emitters, \code{verilog} and \code{vhdl},
with the width inference an extension needs and the hoisting of a
slice of an expression into a wire; the ports sidecar a testbench
generator reads; and, first in the file, the structural skeleton from
the trace walk, which sees fields only.

\lstinputlisting[caption={\code{lib/src/netlist.rs}.},label={lst:r-netlist}]{lib/src/netlist.rs}
